\documentclass[11pt]{article}
\usepackage[margin=1.05in]{geometry}
\usepackage{amsmath,amssymb}
\usepackage{booktabs}
\usepackage[colorlinks=true,linkcolor=blue,urlcolor=blue]{hyperref}
\usepackage{microtype}

\newcommand{\T}{\mathcal{T}}
\newcommand{\A}{\mathcal{A}}
\newcommand{\PP}{\mathcal{P}}
\newcommand{\D}{\mathcal{D}}
\newcommand{\Z}{\mathcal{Z}}
\newcommand{\R}{\mathbb{R}}
\newcommand{\tr}{\mathrm{tr}}

\title{\vspace{-1.5em}Low-rank certification of 2-D PIE stability:\\ the theory behind the demo}
\date{September 2026}
\author{}

\begin{document}
\maketitle
\vspace{-3em}

\begin{abstract}
\noindent
Stability of a 2-D PIE is certified by solving a semidefinite program whose
matrix variables are far larger than the information they carry: the
certificates that exist have rank $2$ per block (and exactly $2n$ for $n$
states), while the matrices being searched over have side length in the
thousands. This note explains, with a minimum of machinery, why that rank is
small, how the demo exploits it by searching only a low-dimensional slice of
the positive semidefinite cone, why a certificate found this way can be
trusted even though the search is nonconvex, and where the method is known to
stop working. Everything quantitative below was measured on the demo's
problem family; nothing is asserted from theory alone.
\end{abstract}

\section{The certification problem, and where the cost lives}

Given a PIE with operators $(\T,\A)$, exponential stability follows from the
existence of a PI operator $\PP\succeq 0$ (coercive) with
\begin{equation}\label{eq:lpi}
\T^{*}\PP\A+\A^{*}\PP\T \;=\; -\,\D, \qquad \D\succeq 0 .
\end{equation}
PIETOOLS turns~\eqref{eq:lpi} into a semidefinite program in the standard
way: $\PP$ and $\D$ are parameterized so that positivity is automatic.
Concretely, a positive PI operator is written as a quadratic form in a fixed
finite basis of simpler operators (monomial multipliers and monomial-kernel
integrals): collecting that basis into a column $\Z(s)$,
\begin{equation}\label{eq:gram}
\PP \;=\; \Z^{*} Q\, \Z, \qquad Q\succeq 0,
\end{equation}
so any positive semidefinite \emph{Gram matrix} $Q$ yields a valid
$\PP\succeq0$, and~\eqref{eq:lpi} becomes a set of \emph{linear} equality
constraints on the entries of the two Gram matrices ($Q_1$ for $\PP$, $Q_2$
for $\D$), obtained by matching the polynomial coefficients of every operator
parameter. The decision problem is therefore
\begin{equation}\label{eq:sdp}
\text{find } Q_1,Q_2\succeq 0 \quad\text{s.t.}\quad
\mathcal{L}(Q_1,Q_2)=b ,
\end{equation}
with $m$ scalar equality constraints and $\sum_i N_i(N_i+1)/2$ scalar
unknowns, where $N_i$ is the side length of Gram block $i$.

The side length is $N_i = (\text{states})\times(\text{basis functions})$, and
in two spatial dimensions the basis is what explodes: a basis over the four
variables $(x,\theta_x,y,\theta_y)$ at per-variable degree $d$ has
$O(d^4)$ elements. At per-variable degree~$4$ the negativity block of the
demo problem already has $N_2=744$ for one state, i.e.\ $277{,}176$ Gram
unknowns; four states give $4.43$ million.

An interior-point solver does not mind the number of unknowns per se; what it
must do at every iteration is form and factor an $m\times m$ system (the
Schur complement of the KKT equations). For these LPIs that matrix is
\emph{completely dense} --- measured on the demo family at $m=35{,}800$:
every one of the $m^2$ entries is structurally nonzero --- so each iteration
costs $O(m^2)$ memory and $O(m^3)$ flops, and $m$ grows like
$(\text{states})^2$. Measured consequences: the one-state problem solves in
$1707$\,s; the four-state problem needs a $66.8$\,GiB dense matrix and cannot
be attempted on a $64$\,GiB machine. Memory, not time, is what fails first.

\section{The observation: certificates have tiny rank}

The matrices being searched over are enormous, but the certificates that
exist inside them are not. Measured on the demo family, with every accepted
certificate verified as described in Section~\ref{sec:verify}:
\begin{itemize}
\item the minimum rank of each Gram block is $2$ for one state, and it stays
      $2$ while the basis (and hence $N$) grows from $424$ to $1208$;
\item across $n=1,\dots,6$ replicated states the minimum rank is exactly
      $2n$ in both blocks --- rank grows with \emph{states}, never with
      \emph{basis size}.
\end{itemize}

There is a mechanism behind this, and it predicts both halves. A positive
operator $\PP\succeq0$ has a square root: $\PP=\mathcal{R}^{*}\mathcal{R}$.
If $\PP$ is an \emph{integral} operator, its square root can itself be
represented in the same operator basis $\Z$, as a small number of rows of
coefficients --- one row per ``square'' in the representation. Writing
$\mathcal{R}$ with $r$ rows of coefficients $Y\in\R^{N\times r}$ gives
$Q=YY^{\top}$, i.e.\ $\mathrm{rank}\,Q\le r$. The point is that $r$ counts
\emph{how many squares the operator needs}, which does not change when the
basis is enriched: raising the polynomial degree, or moving from one spatial
dimension to two, adds columns to the search space without adding squares to
the answer. Adding states does add squares --- the square root of an
$n\times n$ block operator needs $n$ block rows --- which is exactly the
measured $2n$.

Two practical corollaries. First, the useful axes are separated: problem
difficulty (basis size, spatial dimension, degree) inflates $N$ but not the
rank, so the \emph{information content} of a certificate is a handful of
numbers even when the SDP has millions of unknowns. Second, the rank tells
you what a solver \emph{could} exploit; interior-point methods cannot ---
they converge to maximum-rank solutions by design (the returned Gram ranks
were measured at $417$--$423$ out of $424$) --- so the structure is invisible
unless one goes looking for it.

\section{Exploiting it: search a slice, not the cone}

Suppose we know (or guess) an $N\times r$ matrix $V$ with orthonormal columns
such that a certificate exists of the form
\begin{equation}\label{eq:face}
Q \;=\; V S V^{\top}, \qquad S\in\R^{r\times r},\; S\succeq 0 .
\end{equation}
Substituting~\eqref{eq:face} into~\eqref{eq:sdp} leaves a semidefinite
program in $S$ alone: $r(r+1)/2$ unknowns instead of $N(N+1)/2$. With
$r=2$ and $N=744$ that is $3$ unknowns instead of $277{,}140$ for that block; the resulting
SDP solves in seconds. (In convex-analysis language, \eqref{eq:face}
restricts the search to a \emph{face} of the positive semidefinite cone; no
property of faces beyond the substitution itself is used.)

The step that makes this safe is elementary but load-bearing:
\begin{quote}
$S\succeq0$ implies $VSV^{\top}\succeq0$, and the constraints
$\mathcal{L}(Q)=b$ are linear, so \textbf{any} $(S\succeq0)$ satisfying the
restricted constraints yields a genuine solution of the original
problem~\eqref{eq:sdp}. The restriction can only \emph{lose} feasible points;
it can never manufacture one.
\end{quote}
So a bad guess for $V$ costs a failed attempt, never a false certificate.
The converse direction carries no guarantee, and the demo is explicit about
this: a rank at which the search fails is reported as ``not reached,'' never
as a proven lower bound.

\section{Finding the slice --- the weak point, examined}\label{sec:bm}

Everything so far is either measured fact or elementary algebra. The genuine
weakness of the method is concentrated in one place: \emph{$V$ must be
guessed}, and a wrong guess, while never unsafe (Section~3), yields nothing.
This section explains how the demo guesses, why that choice is defensible,
how it fails, and what the alternatives are.

\subsection*{The proposer: factored descent}

The demo finds $V$ by \emph{factored descent} (the Burer--Monteiro
method~\cite{bm2003}): substitute $Q=YY^{\top}$ with $Y\in\R^{N\times r}$
directly into the equality constraints and minimize the residual
$\|\mathcal{L}(YY^{\top})-b\|^2$ over $Y$ by nonlinear least squares
(Levenberg--Marquardt), from a handful of random starting points plus a
warm start from the previous rank. Two things make this the right tool here:
\begin{itemize}
\item the semidefinite constraint disappears --- $YY^{\top}$ is positive
      semidefinite by construction --- leaving an unconstrained problem in
      $N\!\cdot\!r$ variables (a few thousand at $r=2$, not millions), whose
      per-iteration linear algebra involves an $(Nr)\times(Nr)$ system;
\item nothing in it ever touches an $m\times m$ matrix, so the memory wall of
      Section~1 never appears, at any problem size.
\end{itemize}
$Y$ is then orthonormalized to give $V$, the small restricted
SDP~\eqref{eq:face} is solved exactly, and the verification step decides.
The unknown rank is handled by scanning $r$ upward under the same gate, then
shrinking the certified face one direction at a time; overshooting $r$ is
harmless, since a rank-$r$ face contains every lower-rank certificate built
from its directions.

\subsection*{What can be said in its favor --- and what cannot}

The honest theoretical position is asymmetric. In one direction there is a
guarantee of \emph{existence}: any feasible SDP with $m$ equality constraints
admits a solution of rank $r$ with $r(r+1)/2\le m$~\cite{pataki}, so
low-rank solutions ($r\sim\sqrt{2m}$, here 100--270) always exist. In the
other direction there is a guarantee about the \emph{search}: when $r$ is
taken large enough that $r(r+1)/2 > m$, the factored problem has, for
generic data, no spurious local minima --- descent finds a global
solution~\cite{boumal}. The demo operates far below that threshold ($r=2$
against $m$ in the thousands), where no such guarantee applies and stalls in
local minima are possible in principle. What can be claimed instead is
empirical, and it was measured rather than assumed:
\begin{itemize}
\item at the ranks where certificates exist, the descent found them at the
      \emph{first} random seed in most measured cases, and within three seeds
      plus a warm start in the rest;
\item its characteristic failure is a \emph{near-miss}, not divergence: a
      residual settling just above the acceptance gate. One measured case
      stalled at $1.6\times10^{-6}$ at rank $2$ on a problem where a rank-$2$
      certificate provably exists (it was later exhibited by a structural
      construction); the first off-family test showed the same signature.
\end{itemize}
That failure mode is why the architecture is a \emph{proposer--checker}
split, and the split does real work: the nonconvex step only has to find the
right $r$-dimensional \emph{subspace}, not the right point in it --- the
restricted solve then repairs the point exactly within that subspace, and the
% CC, 09/20/2026: "restricted SDP" -> "restricted solve": on the face the
% equality system is determined (np = rank M measured at r = 2..30), so it is
% a linear solve plus a projection onto the PSD cone, not an SDP.  Rank
% deficiency alone is NOT the obstruction: a rank-deficient face (np 21,
% rM 13) is repaired at op rel 1.2e-07 where the SDP call rejects it at
% 2.9e-06.  Severe ill CONDITIONING is: at cond(M) ~ 3e16 neither the
% linear solve nor the SDP finds the certificate the face demonstrably
% contains.
verification step (Section~\ref{sec:verify}) is the sole arbiter. A stalled
descent therefore costs minutes and produces the report ``rank $r$ not
reached,'' which is the correct statement: with this proposer, a failed rank
is never evidence of a rank floor.

For a new operator the discovery step is a one-time cost (measured 11--25
minutes on two cold problems) --- cheaper than a single full interior-point
solve, but not free. Everything after it is.

\subsection*{Alternatives, and their measured status}

Guessing $V$ is a heuristic, so the natural question is what else could
supply it. Each of the following was either measured on this problem class or
is flagged as untested:
\begin{description}
\item[Truncating the interior-point solution.] The obvious first idea, and
  it fails for a structural reason: interior-point methods converge to the
  \emph{maximum}-rank point of the solution set (measured Gram ranks
  417--423 of 424), and the leading eigenvectors of such a point are not the
  low-rank subspace. Faces truncated from full solves certified nothing on
  the demo family --- including in a control experiment where the restriction
  was made generous enough that success was expected.
\item[Trace minimization.] The standard convex surrogate for low rank
  (minimize $\tr\,Q$, then truncate). It supplied usable faces on related
  1-D problems, but on the demo family it was measured to fail: the solver
  exits with a numerical error and the trace objective crushes the small
  Lyapunov block to noise, so its truncated faces fail at every rank tried.
  Worth attempting on a new problem --- it is deterministic --- but not to be
  relied on.
\item[Structure: replication --- not transfer.] When the problem has
  replicated states, Section~6 removes the guess entirely: the $n=1$ face is
  constructed, not searched, at every larger $n$. That is where structural
  reuse ends. Reusing a face across a \emph{physics} change is measured to
  fail and has been abandoned as a goal: on the demo family, the
  $\lambda=2$ face gives an $O(1)$ operator residual already at
  $\lambda=1.9$, while a cold search re-certifies the new problem at the
  \emph{same} rank --- the certificate's rank survives the parameter change,
  its subspace does not. Every new operating point pays its own discovery.
\item[Physics-informed construction.] The mechanism of Section~2 says the
  face \emph{is} the operator square root written in the basis. When a
  candidate Lyapunov operator is known in closed form, $V$ can be written
  down rather than searched: on a 1-D beam example, injecting the known
  energy $\mathrm{diag}(1,c)$ reproduced the certificate to machine
  precision. For PDEs with a known energy functional this is the cheapest
  route of all, and it is the most promising direction for removing the
  guess in general.
\item[Purpose-built low-rank SDP solvers.] Solvers that natively exploit
  low rank exist --- manifold-optimization methods and sketching-based
  conditional-gradient methods among them. On 1-D analogues of these LPIs
  they were measured and did not overtake the pipeline above before the
  problem sizes of interest; they have not been re-evaluated in 2-D, where
  the balance could differ. They are alternatives to the \emph{whole}
  pipeline rather than to the guess of $V$.
\end{description}

\section{Trust the check, not the search}\label{sec:verify}

Because the proposer is nonconvex and the solver is numerical, no certificate
is accepted on the word of either. Acceptance is a direct check of the
original operator identity: reconstruct the PI operators $\PP$ and $\D$ from
the candidate Gram matrices and evaluate
\[
\frac{\max\bigl|\,\T^{*}\PP\A+\A^{*}\PP\T+\D\,\bigr|}
     {\max|\,\T^{*}\PP\A+\A^{*}\PP\T\,|}\;<\;10^{-6}
\]
where the maxima run over the coefficients of \emph{every} parameter of the
resulting operators (all 36 parameter cells of a 2-D PI operator --- checking
a subset was measured to pass garbage), together with positive
semidefiniteness of every Gram block. Both halves are necessary: a point
satisfying the equalities to $10^{-11}$ but indefinite was measured, as was
the reverse. In this sense every certificate the demo reports is
approximate --- exact up to a $10^{-6}$ relative residual --- which is the
same standard the underlying SDP solvers work to.

\section{The structure bonus: states replicate for free}

The Gram basis produced by PIETOOLS is organized in groups (one per kind of
basis operator), with the state index nested \emph{inside} each group. For a
family of $n$ replicated states this has a striking consequence: if $X_1$ is
a certified one-state Gram, then placing $I_n\otimes(\cdot)$ \emph{within
each group} produces an $n$-state certificate with no optimization at all ---
write it down, verify it, done. (The naive global $I_n\otimes X_1$, ignoring
the group structure, is wrong and fails; the demo's \texttt{pielr\_tensor}
performs the correct per-group construction and checks the group layout
against the actual Gram sizes before using it.)

Measured: the $n=1$ certificate replicates through $n=6$ with the residual
unchanged to four digits ($9.291\times10^{-7}$), including sizes where the
direct solve is memory-infeasible. This is why the demo's ``solve'' column is
seconds at every $n$: after one discovery at $n=1$, larger $n$ requires only
assembly and verification.

\section{What it buys, and where it stops}

\begin{center}
\begin{tabular}{@{}rrrrrl@{}}
\toprule
$n$ & Gram unknowns & setup (s) & certify (s) & residual & direct solve\\
\midrule
1 & 277{,}176   & 5.5  & 1.6  & $9.3\times10^{-7}$ & 1707 s\\
2 & 1{,}107{,}952 & 26.4 & 6.2  & $9.3\times10^{-7}$ & $\sim$8.6 h (est.)\\
3 & 2{,}492{,}328 & 68.5 & 16.0 & $9.3\times10^{-7}$ & $\sim$98 h (est.)\\
4 & 4{,}430{,}304 & 278  & 47   & $9.3\times10^{-7}$ & infeasible (66.8 GiB)\\
\bottomrule
\end{tabular}
\end{center}

\noindent
Setup (assembling the LPI) is now the dominant cost: the cone side has been
removed from the scaling path, and what remains is program construction.

Known limits, all measured rather than conjectured:
\begin{enumerate}
\item \textbf{Multipliers cost their sum-of-squares length.} The mechanism
      of Section~2 is about integral operators. If the certificate is
      \emph{forced} to reproduce a multiplication operator $a(s)$ --- as in
      weighted inequalities, or Lyapunov weights for variable-coefficient
      PDEs --- then the corresponding Gram sub-block must represent $a$ as a
      sum of squares, and its rank is exactly the number of squares $a$
      needs. In one spatial variable that number is at most $2$ (every
      nonnegative univariate polynomial is a sum of two squares), so the
      effect is benign. In two variables it is unbounded in degree
      (measured: $a=1+x^2+y^2$ forces rank $3$), and a weight that is
      positive on the domain but \emph{not} a sum of squares makes the
      problem infeasible outright at these settings. Pure PDE stability
      problems dodge all of this structurally: their negativity operator has
      no multiplication part whenever $\T$ and $\A$ have none, which holds
      for every 2-D example shipped with PIETOOLS.
\item \textbf{Low rank is a property of certificates with margin.} As the
      operating point approaches the boundary of what the LPI can certify at
      all, the minimum rank climbs toward full (measured: floor $3$ at half
      the certifiable boundary, $6$ at $0.8$, none found $\le 6$ at $0.95$).
      Borderline certificates are not compressible.
\item \textbf{The measured laws are regularities, not theorems.} Rank $2$,
      the $2n$ law, and the replication property are measurements on the
      reaction--diffusion family plus two PIETOOLS library examples, at the
      direct-form stability LPI with no Positivstellensatz terms. The API
      exists precisely so that other systems can be probed; the first
      off-family datapoint (an anisotropic variant) certified at rank
      $[3\;3]$ with rank $2$ not reached --- consistent with the picture,
      but a reminder that the constants are problem-dependent.
\end{enumerate}

\begin{thebibliography}{9}\small
\bibitem{bm2003} S.~Burer and R.D.C.~Monteiro,
``A nonlinear programming algorithm for solving semidefinite programs via
low-rank factorization,'' \emph{Mathematical Programming} 95(2):329--357,
2003.
\bibitem{pataki} G.~Pataki,
``On the rank of extreme matrices in semidefinite programs and the
multiplicity of optimal eigenvalues,'' \emph{Mathematics of Operations
Research} 23(2):339--358, 1998; see also A.~Barvinok, \emph{Discrete \&
Computational Geometry} 13:189--202, 1995.
\bibitem{boumal} N.~Boumal, V.~Voroninski, and A.S.~Bandeira,
``The non-convex Burer--Monteiro approach works on smooth semidefinite
programs,'' \emph{NeurIPS} 2016; deterministic guarantees in \emph{Comm.
Pure Appl.\ Math.} 73(3):581--608, 2020.
\end{thebibliography}

\end{document}
